\documentclass[11pt,letterpaper]{article}

\usepackage[margin=1in]{geometry}
\usepackage{amsmath,amssymb,amscd,amsthm,mathtools}
\usepackage{array}
\usepackage{booktabs}
\usepackage{hyperref}
\usepackage{microtype}
\usepackage{enumitem}

\numberwithin{equation}{section}

\theoremstyle{plain}
\newtheorem{theorem}{Theorem}[section]
\newtheorem{lemma}[theorem]{Lemma}
\newtheorem{proposition}[theorem]{Proposition}
\newtheorem{corollary}[theorem]{Corollary}
\theoremstyle{definition}
\newtheorem{definition}[theorem]{Definition}
\newtheorem{example}[theorem]{Example}
\theoremstyle{remark}
\newtheorem{remark}[theorem]{Remark}

\newcommand{\NN}{\mathbb{N}}
\newcommand{\RR}{\mathbb{R}}
\newcommand{\ZZ}{\mathbb{Z}}
\newcommand{\EE}{\mathbb{E}}
\newcommand{\PP}{\mathbb{P}}
\newcommand{\Xinf}{X_\infty}
\DeclareMathOperator{\Fold}{Fold}
\DeclareMathOperator{\Law}{Law}

\title{A Discrete Stokes Formula for Fold-Truncation Defects Across Resolutions}
\author{Haobo Ma\\[4pt]\small CHRONOAI PTE.\ LTD., Singapore}
\date{April 2026}

\begin{document}
\maketitle

\begin{abstract}
We study one concrete multiscale phenomenon in the Fold tower: folding does not
commute with naive truncation.  For binary windows of length $n$, the direct
prefix truncation to length $m$ and the fold-aware restriction obtained from the
resolution-$n$ folded output produce different words in general.  We encode this
failure of commutation by an adjacent-scale local defect
$\kappa_{m+1\to m}$ and a tower-level global defect $D_{n\to m}$.  The first
main theorem is a discrete Stokes identity expressing $D_{n\to m}$ as the
telescoping xor-sum of projected local defects across the resolution band
$m,\dots,n-1$.  This yields auditable probability bounds, a Hamming--Lipschitz
budget, and a cocycle formula across intermediate scales.  We also show that
eventual summability of adjacent defects forces consistent inverse-limit lifts.
Finally, we record a conditional Haar-mixing criterion: if the defect process
admits a finite-state additive realization whose nontrivial twisted transfer
matrices all have spectral radius strictly below $1$, then its law converges
exponentially to Haar measure on a finite subgroup of $(\ZZ/2\ZZ)^m$.  The
paper is deliberately confined to this defect theory; it does not import the
sofic machinery or gauge-anomaly extensions of the larger manuscript.
\end{abstract}

\tableofcontents

\section{Introduction}

This paper isolates one obstruction in the resolution tower of Fold maps.  The
core folding paper studies the canonical map from raw binary windows to legal
golden-mean words.  Here we ask a more specific multiscale question: if a word
is first truncated and then folded, how far can the result be from what one
gets by first folding at high resolution and then projecting to low resolution?

The answer is subtle enough to deserve its own paper.  Direct truncation on raw
windows does not commute with folding.  The discrepancy is not noise, however.
It admits a rigid bookkeeping scheme across scales.  At each adjacent step
$m+1\to m$ one has a local defect $\kappa_{m+1\to m}$, and across a whole
resolution band one has a global defect $D_{n\to m}$.  The central result is a
discrete Stokes formula:
\[
\text{global defect}
\;=\;
\text{sum of projected local defects}.
\]
The sum is taken in the xor group $(\ZZ/2\ZZ)^m$, so the identity is exact and
auditable.

The paper makes three contributions.  First, we formulate the local and global
defects and prove the telescoping identity.  Second, we derive quantitative
consequences: an auditable probability bound, a Hamming--Lipschitz budget, and
a cocycle relation across intermediate scales.  Third, we record a conditional
statistical endpoint: once the defect process is realized by a finite-state
additive model with a twisted spectral gap on every nontrivial character, the
defect law converges exponentially to Haar measure on its reachable subgroup.

The point of this paper is therefore narrower than the overall project:
\[
\begin{aligned}
\text{naive truncation}
&\not\!\circ\,\text{Fold} \\
&\Longrightarrow
\text{local defect}
\Longrightarrow
\text{discrete Stokes law}
\Longrightarrow
\text{budget and mixing theory}.
\end{aligned}
\]
We use only the minimal background on Zeckendorf normalization and symbolic
systems; see \cite{Zeckendorf1972,Lekkerkerker1952,Fraenkel1985} for the
representation theory and
\cite{WaltersErgodicTheory,LindMarcus1995,Kitchens1998} for the dynamical
framework.  Mixing estimates follow from standard Markov chain techniques
\cite{LevinPeresWilmer2009MarkovMixing}, using transfer operator methods
surveyed in \cite{ParryPollicott1990,Baladi2000}.  No sofic-image theory,
Fischer covers, or arithmetic anomaly expansions are needed here.

\subsection{Related Work and Contribution}

Zeckendorf representations and Fibonacci numeration have a substantial
literature.  The foundational uniqueness results go back to
\cite{Lekkerkerker1952,Zeckendorf1972}, with broader numeration-system
perspectives in \cite{Fraenkel1985,Grabner1995,BerstelReutenauer2011}.  In that
circle of ideas the main themes are canonical representation, normalization,
and automata-theoretic structure.  The Fold maps used here belong to that
background, but the observable studied in this paper is different: we do not
analyze normalization in isolation, but rather the failure of normalization to
commute with prefix restriction across resolutions.

On the symbolic-dynamics side, the ambient space is the golden-mean shift and
its finite-window factors, so factor maps, finite-state presentations, and
compact-group extensions provide the natural context; see
\cite{WaltersErgodicTheory,LindMarcus1995,Kitchens1998,Boyle1983,KitchensSchmidt1989}.
What appears to be absent from that literature is the specific xor-valued defect
observable considered here.  The main contribution of this note is therefore
not a new normalization algorithm, but an exact identity that isolates this
multiscale obstruction and turns it into auditable deterministic estimates.  In
that sense the statistical discussion in Section~7 should be read as a
conditional appendix to the deterministic theory, not as an independently
verified theorem about the Fold tower itself.

\section{Folding and Naive Truncation}

\subsection{Minimal Setup}

For $m\ge 1$, let
\[
\Omega_m:=\{0,1\}^m
\qquad\text{and}\qquad
X_m:=\{x=(x_1,\dots,x_m)\in\{0,1\}^m:\ x_ix_{i+1}=0\text{ for }1\le i<m\}.
\]
Thus $X_m$ is the set of admissible words in the one-sided golden-mean shift.

For $\omega=\omega_1\cdots\omega_m\in\Omega_m$, define its Fibonacci value by
\[
N(\omega):=\sum_{k=1}^m \omega_k F_{k+1}.
\]
Let $Z(N)$ denote the Zeckendorf expansion of $N$ \cite{Zeckendorf1972}.  We
define the fold map by
\[
\Fold_m(\omega):=\pi_{\infty\to m}(Z(N(\omega)))\in X_m,
\]
where $\pi_{\infty\to m}$ truncates an infinite legal sequence to its first $m$
digits.  Equivalently, $\Fold_m$ is Fibonacci evaluation followed by
Zeckendorf normalization and truncation.

For $n\ge m$ we use two truncation maps:
\[
r_{n\to m}:\Omega_n\to\Omega_m,
\qquad
r_{n\to m}(\omega_1\cdots\omega_n):=\omega_1\cdots\omega_m,
\]
and
\[
\pi_{n\to m}:X_n\to X_m,
\qquad
\pi_{n\to m}(x_1,\dots,x_n):=(x_1,\dots,x_m).
\]
The second map is the natural restriction on already folded words.

\begin{lemma}[Functoriality of restriction]
\label{lem:pi-functorial-golden}
If $n\ge k\ge m$, then
\[
\pi_{n\to m}=\pi_{k\to m}\circ \pi_{n\to k}
\qquad\text{and}\qquad
r_{n\to m}=r_{k\to m}\circ r_{n\to k}.
\]
\end{lemma}

\begin{proof}
Both identities are immediate from the definition of prefix truncation.
\end{proof}

\subsection{Failure of Commutation}

The square formed by $r_{n\to m}$ and $\pi_{n\to m}$ does not commute in
general:
\[
\begin{CD}
\Omega_n @>{\Fold_n}>> X_n \\
@V{r_{n\to m}}VV @VV{\pi_{n\to m}}V \\
\Omega_m @>{\Fold_m}>> X_m.
\end{CD}
\]
The basic obstruction already appears at the smallest scales.

\begin{example}[A noncommuting square at $m=1$]
\label{ex:fold-noncommuting-square}
Take $n=2$, $m=1$, and $\omega=11$.  Then
\[
r_{2\to 1}(11)=1,
\qquad
\Fold_1(r_{2\to 1}(11))=\Fold_1(1)=1.
\]
On the other hand,
\[
N(11)=F_2+F_3=3=F_4,
\]
so the Zeckendorf normal form of $N(11)$ has its first two digits equal to
$00$.  Hence
\[
\Fold_2(11)=00,
\qquad
\pi_{2\to 1}(\Fold_2(11))=0.
\]
Therefore
\[
\Fold_1\circ r_{2\to 1}\neq \pi_{2\to 1}\circ \Fold_2.
\]
\end{example}

Example~\ref{ex:fold-noncommuting-square} is exactly the defect that we
formalize below.

\section{Local and Global Defects}

\subsection{Local Defect}

We work in the xor group structure on $\{0,1\}^m$, written with $\oplus$.
For $v\in\{0,1\}^m$, let
\[
\lvert v\rvert_0:=\#\{i:\ v_i\neq 0\}
\]
denote its Hamming weight.

\begin{definition}[Local defect]
\label{def:fold-local-curvature-defect}
For $m\ge 1$ and $\eta\in\Omega_{m+1}$, define the local defect
\[
\kappa_{m+1\to m}(\eta)
:=
\Fold_m(r_{m+1\to m}(\eta))
\oplus
\pi_{m+1\to m}(\Fold_{m+1}(\eta))
\in \{0,1\}^m.
\]
Its zero/nonzero indicator is
\[
K_{m+1\to m}(\eta):=\mathbf{1}_{\{\kappa_{m+1\to m}(\eta)\neq 0\}}.
\]
\end{definition}

The local defect records the failure of commutation at a single adjacent
resolution step.

\subsection{Global Tower Defect}

\begin{definition}[Global tower defect]
\label{def:fold-global-stokes-defect}
Let $n>m$.  For $\omega\in\Omega_n$, define
\[
D_{n\to m}(\omega)
:=
\Fold_m(r_{n\to m}(\omega))
\oplus
\pi_{n\to m}(\Fold_n(\omega))
\in \{0,1\}^m.
\]
\end{definition}

The defect $D_{n\to m}$ measures the aggregate mismatch accumulated when one
compares the two paths from $\Omega_n$ down to $X_m$ in the noncommuting square.

\begin{table}[ht]
\centering
\small
\begin{tabular}{>{$\displaystyle}l<{$} >{$\displaystyle}l<{$} p{3.8in}}
\toprule
\text{observable} & \text{range} & \text{meaning} \\
\midrule
\kappa_{m+1\to m}(\eta) & \{0,1\}^m &
adjacent-scale local noncommutation defect \\
D_{n\to m}(\omega) & \{0,1\}^m &
aggregate defect across the whole band $m,\dots,n-1$ \\
K_{m+1\to m}(\eta) & \{0,1\} &
indicator that the local defect is nonzero \\
\lvert D_{n\to m}(\omega)\rvert_0 & \{0,\dots,m\} &
Hamming size of the global defect \\
\bottomrule
\end{tabular}
\caption{The defect observables used throughout the paper.}
\label{tab:defect-observables}
\end{table}

\section{A Discrete Stokes Identity}

\subsection{Telescoping Structure}

\begin{theorem}[Discrete Stokes identity for Fold defects]
\label{thm:fold-discrete-stokes-defect}
For every $n>m$ and every $\omega\in\Omega_n$,
\[
\boxed{
D_{n\to m}(\omega)
=
\bigoplus_{j=m}^{n-1}
\pi_{j\to m}\!\left(
\kappa_{j+1\to j}(r_{n\to j+1}(\omega))
\right).
}
\]
\end{theorem}

\begin{proof}
For $j\in\{m,\dots,n\}$ define
\[
A_j(\omega):=\Fold_j(r_{n\to j}(\omega)),
\qquad
B_j(\omega):=\pi_{n\to j}(\Fold_n(\omega)),
\]
and
\[
\Delta_j(\omega):=A_j(\omega)\oplus B_j(\omega)\in\{0,1\}^j.
\]
Then $\Delta_m(\omega)=D_{n\to m}(\omega)$ and $\Delta_n(\omega)=0$.

Fix $j\in\{m+1,\dots,n\}$.  By Lemma~\ref{lem:pi-functorial-golden},
\[
B_{j-1}(\omega)=\pi_{j\to j-1}(B_j(\omega))
\]
and
\[
A_{j-1}(\omega)=\Fold_{j-1}(r_{j\to j-1}(r_{n\to j}(\omega))).
\]
Therefore
\[
\begin{aligned}
\Delta_{j-1}(\omega)
&=A_{j-1}(\omega)\oplus B_{j-1}(\omega) \\
&=\Bigl(A_{j-1}(\omega)\oplus \pi_{j\to j-1}(A_j(\omega))\Bigr)
\oplus
\pi_{j\to j-1}\bigl(A_j(\omega)\oplus B_j(\omega)\bigr).
\end{aligned}
\]
The first bracket is exactly
\[
\kappa_{j\to j-1}(r_{n\to j}(\omega)),
\]
while the second term is $\pi_{j\to j-1}(\Delta_j(\omega))$.  Hence
\[
\Delta_{j-1}(\omega)
=
\kappa_{j\to j-1}(r_{n\to j}(\omega))
\oplus
\pi_{j\to j-1}(\Delta_j(\omega)).
\]
Iterating from $j=n$ down to $j=m+1$ and using the fact that prefix projection
commutes with xor gives
\[
\Delta_m(\omega)
=
\bigoplus_{j=m}^{n-1}
\pi_{j\to m}\!\left(
\kappa_{j+1\to j}(r_{n\to j+1}(\omega))
\right).
\]
Since $\Delta_m(\omega)=D_{n\to m}(\omega)$, the proof is complete.
\end{proof}

\begin{remark}[Why ``Stokes''?]
The terminology is analogical rather than literal.  The quantity
$D_{n\to m}(\omega)$ is a boundary-level discrepancy between two routes from
resolution $n$ to resolution $m$, while the adjacent defects
$\kappa_{j+1\to j}$ are the interior contributions picked up one scale at a
time.  Theorem~\ref{thm:fold-discrete-stokes-defect} therefore has the same
formal shape as a discrete boundary-equals-sum-of-local-contributions identity.
It is closer to a coboundary computation in a finite cochain complex than to a
differential-geometric Stokes theorem; compare
\cite{Hatcher2002,Forman1998}.
\end{remark}

\begin{example}[A worked Stokes computation at $n=3$, $m=1$]
\label{ex:fold-stokes-computation}
Take $\omega=111$.  Then
\[
r_{3\to 1}(111)=1,
\qquad
\Fold_1(r_{3\to 1}(111))=1.
\]
Also
\[
N(111)=F_2+F_3+F_4=6=F_2+F_5,
\]
so the Zeckendorf expansion begins with the digits $1001\cdots$, and hence
\[
\Fold_3(111)=100,
\qquad
\pi_{3\to 1}(\Fold_3(111))=1.
\]
Therefore the global defect vanishes:
\[
D_{3\to 1}(111)=1\oplus 1=0.
\]
At the adjacent scales we have
\[
\kappa_{2\to 1}(r_{3\to 2}(111))
=
\kappa_{2\to 1}(11)
=
\Fold_1(1)\oplus \pi_{2\to 1}(\Fold_2(11))
=
1\oplus 0
=
1,
\]
while
\[
\kappa_{3\to 2}(111)
=
\Fold_2(11)\oplus \pi_{3\to 2}(\Fold_3(111))
=
00\oplus 10
=
10.
\]
Projecting the second term to level $m=1$ gives $\pi_{2\to 1}(10)=1$, so
\[
\kappa_{2\to 1}(11)\oplus \pi_{2\to 1}(\kappa_{3\to 2}(111))
=
1\oplus 1
=
0
=
D_{3\to 1}(111).
\]
This is exactly the discrete Stokes identity in the first nontrivial resolution
band where two adjacent defects contribute.
\end{example}

\subsection{Auditable Consequences}

\begin{corollary}[Auditable probability bound]
\label{cor:fold-discrete-stokes-auditable-bound}
Let $n>m$, let $\mu$ be any probability measure on $\Omega_n$, and let
$f:X_m\to\RR$ be bounded.  Then
\[
\left|
\EE_\mu f\bigl(\Fold_m(r_{n\to m}(\omega))\bigr)
-
\EE_\mu f\bigl(\pi_{n\to m}(\Fold_n(\omega))\bigr)
\right|
\le
2\lVert f\rVert_\infty\,
\PP_\mu[D_{n\to m}(\omega)\neq 0].
\]
Moreover,
\[
\{D_{n\to m}\neq 0\}
\subseteq
\bigcup_{j=m}^{n-1}
\left\{
K_{j+1\to j}(r_{n\to j+1}(\omega))=1
\right\},
\]
and therefore
\[
\PP_\mu[D_{n\to m}\neq 0]
\le
\sum_{j=m}^{n-1}
\PP_\mu\!\left[
K_{j+1\to j}(r_{n\to j+1}(\omega))=1
\right].
\]
\end{corollary}

\begin{proof}
Write
\[
x=\Fold_m(r_{n\to m}(\omega)),
\qquad
y=\pi_{n\to m}(\Fold_n(\omega)).
\]
Then $\{x\neq y\}=\{D_{n\to m}(\omega)\neq 0\}$, and
\[
\lvert f(x)-f(y)\rvert
\le
2\lVert f\rVert_\infty\,\mathbf{1}_{\{x\neq y\}}.
\]
Taking expectation yields the first inequality.

If every local defect in the Stokes sum of
Theorem~\ref{thm:fold-discrete-stokes-defect} vanishes, then the xor-sum vanishes as
well.  The event inclusion follows by contraposition, and the final estimate is
just the union bound.
\end{proof}

\section{Budget Bounds and Cocycle Structure}

\subsection{Hamming--Lipschitz Budget}

\begin{corollary}[Hamming--Lipschitz budget]
\label{cor:fold-hamming-lipschitz-budget}
Let $n>m$, let $\mu$ be any probability measure on $\Omega_n$, and let
$d_{\mathrm H}(x,y):=\lvert x\oplus y\rvert_0$ be Hamming distance on $X_m$.
If $f:X_m\to\RR$ is $L$-Lipschitz with respect to $d_{\mathrm H}$, then
\[
\left|
\EE_\mu f\bigl(\Fold_m(r_{n\to m}(\omega))\bigr)
-
\EE_\mu f\bigl(\pi_{n\to m}(\Fold_n(\omega))\bigr)
\right|
\le
L\,\EE_\mu \lvert D_{n\to m}(\omega)\rvert_0
\]
and
\[
\EE_\mu \lvert D_{n\to m}(\omega)\rvert_0
\le
\sum_{j=m}^{n-1}
\EE_\mu\left\lvert
\kappa_{j+1\to j}(r_{n\to j+1}(\omega))
\right\rvert_0.
\]
\end{corollary}

\begin{proof}
For the same $x$ and $y$ as in the proof of
Corollary~\ref{cor:fold-discrete-stokes-auditable-bound},
\[
\lvert f(x)-f(y)\rvert
\le
L\,d_{\mathrm H}(x,y)
=
L\,\lvert D_{n\to m}(\omega)\rvert_0.
\]
Taking expectation gives the first inequality.

For the second, apply Theorem~\ref{thm:fold-discrete-stokes-defect} and use the
subadditivity
\[
\lvert u\oplus v\rvert_0\le \lvert u\rvert_0+\lvert v\rvert_0
\]
together with the fact that prefix projection cannot increase Hamming weight.
\end{proof}

\subsection{Cocycle and Triangle Structure}

\begin{proposition}[Cocycle identity across intermediate scales]
\label{prop:fold-defect-cocycle}
If $m<k<n$, then for every $\omega\in\Omega_n$,
\[
D_{n\to m}(\omega)
=
D_{k\to m}(r_{n\to k}(\omega))
\oplus
\pi_{k\to m}(D_{n\to k}(\omega)).
\]
Consequently, if
\[
\Delta_{n,m}(\omega):=\lvert D_{n\to m}(\omega)\rvert_0,
\]
then
\[
\Delta_{n,m}(\omega)
\le
\Delta_{k,m}(r_{n\to k}(\omega))
+
\Delta_{n,k}(\omega).
\]
\end{proposition}

\begin{proof}
Split the Stokes sum of Theorem~\ref{thm:fold-discrete-stokes-defect} at the
intermediate scale $k$:
\[
\begin{aligned}
D_{n\to m}(\omega)
&=
\bigoplus_{j=m}^{k-1}
\pi_{j\to m}\!\left(
\kappa_{j+1\to j}(r_{n\to j+1}(\omega))
\right) \\
&\quad\oplus
\bigoplus_{j=k}^{n-1}
\pi_{j\to m}\!\left(
\kappa_{j+1\to j}(r_{n\to j+1}(\omega))
\right).
\end{aligned}
\]
The first block depends only on the prefix $r_{n\to k}(\omega)$ and is exactly
$D_{k\to m}(r_{n\to k}(\omega))$.  In the second block we use
$\pi_{j\to m}=\pi_{k\to m}\circ \pi_{j\to k}$ from
Lemma~\ref{lem:pi-functorial-golden} and factor out $\pi_{k\to m}$, obtaining
$\pi_{k\to m}(D_{n\to k}(\omega))$.  This proves the cocycle identity.

The triangle inequality follows from
\[
\lvert u\oplus v\rvert_0\le \lvert u\rvert_0+\lvert v\rvert_0
\]
and the nonexpansiveness of prefix projection in Hamming weight.
\end{proof}

\section{Lifts and Asymptotic Compatibility}

\subsection{Consistent Lifts}

Let
\[
\Xinf:=\{x=(x_1,x_2,\dots)\in\{0,1\}^{\NN}: x_ix_{i+1}=0\text{ for all }i\ge 1\}
\]
be the one-sided golden-mean shift.  For a raw infinite binary sequence
$\omega^\infty=(\omega_1,\omega_2,\dots)\in\{0,1\}^{\NN}$, let
\[
r_{\infty\to m}(\omega^\infty):=(\omega_1,\dots,\omega_m)\in\Omega_m,
\]
and for $x\in\Xinf$ let
\[
\pi_{\infty\to m}(x):=(x_1,\dots,x_m)\in X_m.
\]

\begin{proposition}[Consistent lift after eventual defect extinction]
\label{prop:fold-naive-prefix-lift}
Let $\omega^\infty\in\{0,1\}^{\NN}$ and define
\[
x_m:=\Fold_m(r_{\infty\to m}(\omega^\infty))\in X_m.
\]
For each $m\ge 1$, define the adjacent discrepancy
\[
\delta_m:=x_m\oplus \pi_{m+1\to m}(x_{m+1})\in\{0,1\}^m.
\]
Note that
\[
\delta_m=\kappa_{m+1\to m}(r_{\infty\to m+1}(\omega^\infty)),
\]
so the summability hypothesis is exactly the summability of the adjacent local
defects along the Fold trajectory of $\omega^\infty$.
If
\[
\sum_{m\ge 1}\lvert \delta_m\rvert_0<\infty,
\]
then there exists $M$ such that $\delta_m=0$ for every $m\ge M$.  Hence
$(x_m)_{m\ge M}$ is compatible under restriction, and there exists a unique
$x^\infty\in\Xinf$ satisfying
\[
\pi_{\infty\to m}(x^\infty)=x_m
\qquad\text{for all }m\ge M.
\]
\end{proposition}

\begin{proof}
Since each $\lvert\delta_m\rvert_0$ is a nonnegative integer and their sum is
finite, only finitely many of them can be nonzero.  Therefore there exists $M$
such that $\delta_m=0$ for all $m\ge M$, which means
\[
x_m=\pi_{m+1\to m}(x_{m+1})
\qquad\text{for every }m\ge M.
\]
Thus $(x_m)_{m\ge M}$ is a compatible family of admissible prefixes, so it
defines a unique infinite admissible sequence $x^\infty\in\Xinf$.
\end{proof}

\begin{remark}
Proposition~\ref{prop:fold-naive-prefix-lift} is the precise form of asymptotic
compatibility needed for this paper.  We do not claim that all Fold trajectories
stabilize in this way.  The proposition only says that summable adjacent defects
force eventual consistency.
\end{remark}

\section{Statistical Remarks}

\subsection{A Conditional Haar-Mixing Criterion}

The next theorem is a conditional criterion.  It isolates the exact ingredients
needed to deduce Haar mixing once the defect process has been encoded by a
finite-state additive model.

\begin{theorem}[Conditional Haar mixing from a finite-state additive model]
\label{thm:fold-stokes-defect-haar-mixing}
Fix $m\ge 1$, and let $\omega$ be uniformly distributed on $\Omega_n$ for
$n\ge m$.  Suppose there exist a finite state space $\mathsf{S}_m$, a
primitive stochastic matrix $P_m=(p_{ab})_{a,b\in\mathsf{S}_m}$, an initial
distribution $\nu_m$ on $\mathsf{S}_m$, a finite subgroup
$\mathsf{H}_m\le(\ZZ/2\ZZ)^m$, and an increment map
$\Gamma_m:\mathsf{S}_m\times\mathsf{S}_m\to\mathsf{H}_m$ such that:
\begin{enumerate}[label=(\roman*)]
\item if $(Y_t)_{t\ge 0}$ is the Markov chain with initial law $\nu_m$ and
transition matrix $P_m$, then for every $n\ge m$ the defect satisfies the
distributional identity
\[
D_{n\to m}
\stackrel{\mathrm{law}}=
\bigoplus_{t=0}^{n-m-1}\Gamma_m(Y_t,Y_{t+1});
\]
\item for every nontrivial character
$\chi\in\widehat{\mathsf{H}_m}\setminus\{1\}$, the twisted matrix
\[
P_{m,\chi}(a,b):=p_{ab}\,\chi(\Gamma_m(a,b))
\qquad (a,b\in\mathsf{S}_m)
\]
satisfies $\rho(P_{m,\chi})<1$.
\end{enumerate}
Then there exist constants $C_m>0$ and $0<\rho_m<1$ such that for all $n\ge m$,
\[
\left\|
\Law(D_{n\to m})-\mathrm{Haar}(\mathsf{H}_m)
\right\|_{\mathrm{TV}}
\le
C_m\,\rho_m^{\,n-m}.
\]
\end{theorem}

\begin{remark}[Standing hypotheses]
\label{rem:standing-hypotheses-haar-mixing}
Assumption (ii) is the genuinely conditional input.  In a concrete Fold model
one would derive $\mathsf{S}_m$, $P_m$, and $\Gamma_m$ from a finite transducer
and then verify $\rho(P_{m,\chi})<1$ for each nontrivial character by direct
matrix computation.  The present paper does not carry out that verification, so
Section~7 should be read as a conditional criterion rather than as an
established property of the Fold tower.  By contrast, the deterministic Stokes
identity (Theorem~\ref{thm:fold-discrete-stokes-defect}) and its corollaries
hold unconditionally.
\end{remark}

\begin{proof}[Proof of Theorem~\ref{thm:fold-stokes-defect-haar-mixing}]
Set $\ell:=n-m$ and let $\mu_\ell:=\Law(D_{n\to m})$, viewed as a probability
measure on $\mathsf{H}_m$.  For $\chi\in\widehat{\mathsf{H}_m}$, write
\[
\widehat{\mu_\ell}(\chi):=\sum_{h\in\mathsf{H}_m}\mu_\ell(h)\chi(h)
=
\EE[\chi(D_{n\to m})].
\]
Using assumption (i) and expanding over all state paths of the Markov chain,
we obtain
\[
\widehat{\mu_\ell}(\chi)
=
\nu_m P_{m,\chi}^{\,\ell}\mathbf{1},
\]
where $\mathbf{1}$ denotes the all-ones column vector on $\mathsf{S}_m$.

For the trivial character $\chi=1$, this gives
$\widehat{\mu_\ell}(1)=\nu_m P_m^{\,\ell}\mathbf{1}=1$.  For
$\chi\neq 1$, assumption (ii) gives $\rho(P_{m,\chi})<1$.  Choose
$\rho_m$ such that
\[
\max_{\chi\neq 1}\rho(P_{m,\chi})<\rho_m<1.
\]
Since there are only finitely many nontrivial characters and each $P_{m,\chi}$
is a finite matrix, standard finite-dimensional spectral theory yields
constants $C_\chi>0$ such that
\[
\lVert P_{m,\chi}^{\,\ell}\rVert\le C_\chi\,\rho_m^{\,\ell}
\qquad\text{for all }\ell\ge 0.
\]
Hence
\[
\lvert\widehat{\mu_\ell}(\chi)\rvert
\le
\lVert \nu_m\rVert\,
\lVert P_{m,\chi}^{\,\ell}\rVert\,
\lVert \mathbf{1}\rVert
\le
C'_\chi\,\rho_m^{\,\ell}
\]
for suitable constants $C'_\chi$.

Now Fourier inversion on the finite abelian group $\mathsf{H}_m$ gives, for
each $h\in\mathsf{H}_m$,
\[
\mu_\ell(h)-\frac{1}{|\mathsf{H}_m|}
=
\frac{1}{|\mathsf{H}_m|}
\sum_{\chi\neq 1}\widehat{\mu_\ell}(\chi)\,\overline{\chi(h)}.
\]
Therefore
\[
\begin{aligned}
\left\|
\mu_\ell-\mathrm{Haar}(\mathsf{H}_m)
\right\|_{\mathrm{TV}}
&=
\frac12\sum_{h\in\mathsf{H}_m}
\left|
\mu_\ell(h)-\frac{1}{|\mathsf{H}_m|}
\right| \\
&\le
\frac12\sum_{\chi\neq 1}\lvert\widehat{\mu_\ell}(\chi)\rvert \\
&\le
\frac12\sum_{\chi\neq 1}C'_\chi\,\rho_m^{\,\ell}.
\end{aligned}
\]
Setting $C_m:=\frac12\sum_{\chi\neq 1}C'_\chi$ and recalling that
$\ell=n-m$ proves the stated estimate.
\end{proof}

\subsection{Parity Polarization}

\begin{corollary}[Parity polarization and binomial limit at full generation]
\label{cor:fold-stokes-defect-parity-polarization}
Under the hypotheses of Theorem~\ref{thm:fold-stokes-defect-haar-mixing}, fix
$u\in (\ZZ/2\ZZ)^m$ and define
\[
\chi_u(v):=(-1)^{\langle u,v\rangle}
\qquad (v\in (\ZZ/2\ZZ)^m).
\]
Then there exist constants $C_m>0$ and $0<\rho_m<1$ such that for all $n\ge m$,
\[
\left|
\EE[\chi_u(D_{n\to m})]
-
\mathbf{1}_{\{u\in \mathsf{H}_m^\perp\}}
\right|
\le
C_m\,\rho_m^{\,n-m},
\]
where
\[
\mathsf{H}_m^\perp:=\{u\in(\ZZ/2\ZZ)^m:\chi_u|_{\mathsf{H}_m}\equiv 1\}.
\]
In particular, if $\mathsf{H}_m=(\ZZ/2\ZZ)^m$, then
\[
\EE[\chi_u(D_{n\to m})]=O(\rho_m^{\,n-m})
\qquad\text{for every }u\neq 0,
\]
and
\[
\lvert D_{n\to m}\rvert_0
\xrightarrow[n\to\infty]{\mathrm{law}}
\mathrm{Binomial}\!\left(m,\frac12\right).
\]
\end{corollary}

\begin{proof}
The total-variation estimate from
Theorem~\ref{thm:fold-stokes-defect-haar-mixing} controls the expectation of
every $\{\pm 1\}$-valued observable, hence of $\chi_u$.  Under Haar measure on
$\mathsf{H}_m$, the character average equals $1$ exactly on the annihilator
$\mathsf{H}_m^\perp$ and vanishes otherwise.  This gives the first claim.

If $\mathsf{H}_m=(\ZZ/2\ZZ)^m$, the limit law is uniform on all binary vectors
of length $m$.  Under that law the coordinates are independent Bernoulli$(1/2)$
variables, so the Hamming weight has the stated binomial distribution.
\end{proof}

\begin{remark}
The statistical section is intentionally kept short.  The main paper here is
the deterministic defect identity.  The mixing theorem is included because it
shows that the defect process is not merely a bookkeeping device but a genuine
dynamical observable.
\end{remark}

\section{Discussion}

The paper isolates one exact mechanism: folding and naive truncation fail to
commute, but the failure is structured.  Local adjacent-scale defects assemble
into a global tower defect by a discrete Stokes formula.  Once this identity is
available, quantitative corollaries follow with little extra work: auditable
probability bounds, Lipschitz budgets, and cocycle inequalities.

The consistent-lift proposition gives the clean finite-to-infinite endpoint of
the deterministic theory.  The Haar-mixing theorem shows how the same defect can
be treated statistically when the Fold tower is read by a finite-state additive
kernel.  This is the correct stopping point for the present note.  We do not
import the broader sofic, gauge-anomaly, or arithmetic expansions of the larger
project.

\bibliographystyle{alpha}
\bibliography{references}

\end{document}
